% Control predictivo con horizonte móvil: planear varias decisiones, ejecutar solo la primera.
\begin{tikzpicture}[x=1.45cm,y=1cm,
  plan/.style={circle,draw=qaqua,fill=white,line width=0.8pt,minimum size=2.6mm,inner sep=0pt},
  ejec/.style={circle,draw=qaqua,fill=qaqua,line width=0.8pt,minimum size=3.2mm,inner sep=0pt}]
\node[titulo,anchor=east] at (-0.35,1.5) {Plan en $t$};
\node[titulo,anchor=east] at (-0.35,0) {Plan en $t+1$};
\foreach \k/\y in {0/1.75,1/1.45,2/1.3,3/1.25,4/1.2} \node[plan] (a\k) at (\k,\y) {};
\foreach \k in {0,...,3}{\pgfmathtruncatemacro{\n}{\k+1}\draw[qaqua!60,line width=0.6pt] (a\k) -- (a\n);}
\node[ejec] at (0,1.75) {};
\foreach \k/\y in {1/0.2,2/-0.05,3/-0.15,4/-0.2,5/-0.22} \node[plan] (b\k) at (\k,\y) {};
\foreach \k in {1,...,4}{\pgfmathtruncatemacro{\n}{\k+1}\draw[qaqua!60,line width=0.6pt] (b\k) -- (b\n);}
\node[ejec] at (1,0.2) {};
\draw[guia,line width=0.7pt] (-0.2,-0.75) -- (5.4,-0.75);
\foreach \k/\e in {0/t,1/t+1,2/t+2,3/t+3,4/t+4,5/t+5}{\draw[guia] (\k,-0.75) -- (\k,-0.83);
  \node[nota,anchor=north] at (\k,-0.85) {$\e$};}
\draw[flecha,draw=qnaranja] (0.35,1.05) .. controls (0.6,0.8) .. (0.85,0.5);
\node[nota,anchor=west,text=tinta] at (0.72,0.95) {llegan datos nuevos: se recalcula todo el plan};
\node[nota,anchor=west] at (5.55,1.2) {solo se ejecuta\\ el primer paso\\ (círculo lleno)};
\end{tikzpicture}
